\chapter{Markov decision processes}
\label{ch:mdp}

This chapter gives the mathematical framework of reinforcement learning:
Markov decision processes (MDPs), returns, policies and value functions. Its
central results are the Bellman equations, which characterise the value of
a policy and the value of an optimal policy as fixed points of contractions;
every algorithm of this course is built on them.

Throughout the theory we assume, as in Pacman, that the sets of states and
actions are finite, and that the rewards are bounded.

\section{Markov decision processes}

\begin{definition}{Markov decision process}{mdp}
  A (finite) \emph{Markov decision process} is a tuple
  $(\States, \Actions, p, \gamma)$ where
  \begin{itemize}[nosep]
    \item $\States$ is a finite set of \emph{states} and $\Actions$ a finite
      set of \emph{actions};
    \item $p(s', r \mid s, a)$ is the probability that the environment
      moves to state $s'$ and gives reward $r$ when action $a$ is taken in
      state $s$, the rewards belonging to a finite set $\Rewards \subset
      \R$:
      \begin{equation}
        \sum_{s' \in \States} \sum_{r \in \Rewards} p(s', r \mid s, a) = 1
        \qquad \text{for all } s \in \States,\ a \in \Actions;
      \end{equation}
    \item $\gamma \in [0, 1)$ is the \emph{discount factor}.
  \end{itemize}
  The process generates random variables $S_0, A_0, R_1, S_1, A_1, R_2,
  \dots$ such that, for every $t$ and every history,
  \begin{equation}
    \label{eq:markov}
    \Prob\left(S_{t+1} = s', R_{t+1} = r \given S_t = s, A_t = a,
      H_{t-1}\right) = p(s', r \mid s, a),
  \end{equation}
  where $H_{t-1} = (S_0, A_0, R_1, \dots, S_{t-1}, A_{t-1}, R_t)$ is the
  past. Equation \eqref{eq:markov} is the \emph{Markov property}: the future
  depends on the past only through the current state and action.
\end{definition}

\begin{intuition}
  The state is a summary of the past that is enough to predict the future,
  like a chess position: how the pieces got there does not matter for what
  can happen next. The Markov property is a property of the \emph{choice} of
  state as much as of the system: a process that is not Markov in some
  variables becomes Markov when the state includes enough of the history
  (the direction of the ghosts in game \code{v2}, \cref{rem:pomdp}).
\end{intuition}

Two quantities derived from $p$ are used constantly: the \emph{transition
probabilities} and the \emph{expected rewards}
\begin{equation}
  p(s' \mid s, a) = \sum_{r} p(s', r \mid s, a), \qquad
  r(s, a) = \E\left[R_{t+1} \given S_t = s, A_t = a\right]
  = \sum_{s', r} r\, p(s', r \mid s, a).
\end{equation}

\paragraph{Episodes.} A game ends in a \emph{terminal} state. It is
convenient to model a terminal state as an \emph{absorbing} state: every
action keeps the process in it with reward $0$. With this convention all
the trajectories are infinite, and the episodic and continuing settings are
handled by the same equations.

\begin{pacman}{The game as an MDP}{pacman-mdp}
  A state of game \code{v1} is the configuration of the board: Pacman's
  cell, the cells of the ghosts still in the game, the set of remaining
  cookies, the set of remaining power cookies and Pacman's power credit
  $c \in \{0, \dots, 12\}$; plus two absorbing states, ``won'' and
  ``lost''. The actions are $\Actions = \{\code{E}, \code{N}, \code{W},
  \code{S}, \code{I}\}$. The randomness of $p$ comes only from the ghosts
  (\cref{sec:game}): given the configuration and Pacman's action, the
  probability of each new configuration is a product of the probabilities
  of the moves of the ghosts. The reward is a deterministic function of the
  configuration, the action and the new configuration
  (\cref{tab:rewards}).

  The Markov property holds because the ghosts' moves depend only on the
  current positions and on Pacman's power (\cref{rem:ghost-heading}): the
  past adds nothing once the configuration is known.
\end{pacman}

\begin{remark}{Observations and partial observability}{pomdp}
  The agent does not need the configuration in its Python form: the
  feature planes of \cref{sec:observation} encode exactly the same
  information (walls, cookies, power cookies, Pacman, each ghost, the power
  credit). In game \code{v1} the image is therefore a complete state, and the
  second stacked frame is redundant in principle: it would be necessary if
  the ghosts' moves depended on their direction, which the second frame
  reveals.

  Game \code{v2} differs in two ways. Its ghosts prefer to keep their
  direction (\cref{rem:ghost-heading}), so the direction of each ghost is
  part of the state: the second frame is now needed, and it is enough,
  since the displacement of a ghost between the two frames is its last
  move. Its eaten ghosts also come back. In a first version they
  reappeared on a random cell $15$ moves after being eaten: the time since
  a ghost was eaten then decided when it came back, and it appears in no
  plane of the image. The process seen by the agent was a \emph{partially
  observable} MDP (POMDP): the image was not a state, and stacking frames
  does not help (the timer can be older than two frames). In game
  \code{v2} the eaten ghosts instead wait in a ghost house, visible on the
  board, and once Pacman's power has ended each of them leaves it with
  probability $p = 0.1$ at each move. The waiting time $T$ is geometric,
  $\Prob(T > k) = (1 - p)^k$, hence \emph{memoryless}:
  \begin{equation*}
    \Prob(T > k + m \mid T > k) = \frac{(1 - p)^{k + m}}{(1 - p)^k} = (1 -
    p)^m = \Prob(T > m).
  \end{equation*}
  How long a ghost has already waited tells nothing about when it will
  leave, and the image (with its power plane) is again a state.
\end{remark}

\section{Returns}

\begin{definition}{Return}{return}
  The \emph{return} from time $t$ is the discounted sum of the future
  rewards
  \begin{equation}
    \label{eq:return}
    G_t = R_{t+1} + \gamma R_{t+2} + \gamma^2 R_{t+3} + \cdots
    = \sum_{k=0}^{\infty} \gamma^k R_{t+k+1}.
  \end{equation}
\end{definition}

\begin{proposition}{Properties of the return}{return-props}
  Let $R_{\max} = \max_{r \in \Rewards} \abs{r}$. Then
  \begin{enumerate}[nosep, label=(\roman*)]
    \item the series \eqref{eq:return} converges absolutely and
      $\abs{G_t} \leq R_{\max} / (1 - \gamma)$;
    \item $G_t = R_{t+1} + \gamma\, G_{t+1}$.
  \end{enumerate}
\end{proposition}

\begin{proof}
  (i) $\sum_k \gamma^k \abs{R_{t+k+1}} \leq R_{\max} \sum_k \gamma^k =
  R_{\max} / (1 - \gamma)$. (ii) Since the series converges absolutely we
  may split off its first term:
  $G_t = R_{t+1} + \sum_{k \geq 1} \gamma^k R_{t+k+1}
  = R_{t+1} + \gamma \sum_{j \geq 0} \gamma^j R_{t+1+j+1}
  = R_{t+1} + \gamma\, G_{t+1}$.
\end{proof}

\begin{intuition}
  The return is the total reward from now on, where a reward $k$ steps
  ahead counts only $\gamma^k$ of its value. A useful reading: if at each
  step the game could stop for an outside reason with probability $1 -
  \gamma$, the expected total of the rewards actually collected would be
  exactly $\E[G_t]$; discounting is uncertainty about the future, or
  impatience. Property (ii) is the seed of everything that follows: the
  future seen from now is the next reward plus the (discounted) future seen
  from the next step.
\end{intuition}

The discount factor has two roles. Mathematically it makes the returns
finite; practically it sets a \emph{horizon}: a reward $k$ steps ahead
counts $\gamma^k$, and $\sum_k \gamma^k = 1 / (1 - \gamma)$ is the
effective number of steps the agent looks ahead.

\begin{pacman}{Choosing $\gamma$}{gamma}
  The agents use $\gamma = 0.99$, an effective horizon of $100$ moves,
  about the length of a game. The reward of winning ($+30$) obtained $80$
  moves later is worth $30 \times 0.99^{80} \approx 13.4$ now: still enough
  to make the agent prefer finishing the game over wandering, but less than
  eating a ghost now ($+20$). With $\gamma = 0.9$ (horizon $10$ moves) the
  win would be worth $30 \times 0.9^{80} \approx 0.007$: the agent would
  be blind to it.
\end{pacman}

\begin{remark}{Termination and truncation}{truncation}
  A game \emph{aborted} after the maximal number of moves is not a terminal
  state of the game: the configuration reached still has a value, and the
  move counter, which caused the end, is not part of the state. Treating an
  aborted game as terminal would teach the agent that the states reached
  late in long games have value $0$. All the algorithms of the repository
  therefore distinguish the two cases: after a real end, the return
  stops; after an abort, the return is completed with an estimate of the
  value of the last state (\emph{bootstrapping}, \cref{ch:td}).
\end{remark}

\section{Policies and value functions}

\begin{definition}{Policy}{policy}
  A (stationary, Markov) \emph{policy} is a function $\pi$ that gives, for
  every state $s$, a probability distribution $\pi(\cdot \mid s)$ over
  $\Actions$; the agent following $\pi$ chooses $A_t \sim \pi(\cdot \mid
  S_t)$. The policy is \emph{deterministic} if each $\pi(\cdot \mid s)$ is
  concentrated on one action, which we then write $\pi(s)$.
\end{definition}

Given an initial state, a policy and the MDP define the distribution of the
whole trajectory; we write $\E_\pi$ for expectations under it.

\begin{definition}{Value functions}{values}
  The \emph{state-value function} and the \emph{action-value function} of a
  policy $\pi$ are
  \begin{equation}
    v_\pi(s) = \E_\pi\left[G_t \given S_t = s\right], \qquad
    q_\pi(s, a) = \E_\pi\left[G_t \given S_t = s, A_t = a\right].
  \end{equation}
\end{definition}

By the Markov property these expectations do not depend on $t$ (the
distribution of the future given $S_t = s$ is the same at all times). The
value $v_\pi(s)$ answers ``how good is it to be in $s$ if I follow
$\pi$?''; $q_\pi(s, a)$ answers ``how good is it to play $a$ in $s$ and then
follow $\pi$?''.

\begin{proposition}{Relations between $v_\pi$ and $q_\pi$}{v-q}
  For all states $s$ and actions $a$,
  \begin{align}
    \label{eq:v-from-q}
    v_\pi(s) &= \sum_{a} \pi(a \mid s)\, q_\pi(s, a), \\
    \label{eq:q-from-v}
    q_\pi(s, a) &= \sum_{s', r} p(s', r \mid s, a)
      \left[ r + \gamma\, v_\pi(s') \right]
      = r(s, a) + \gamma \sum_{s'} p(s' \mid s, a)\, v_\pi(s').
  \end{align}
\end{proposition}

\begin{proof}
  \eqref{eq:v-from-q} is the law of total expectation, conditioning on
  $A_t$, whose distribution given $S_t = s$ is $\pi(\cdot \mid s)$. For
  \eqref{eq:q-from-v}, write $G_t = R_{t+1} + \gamma G_{t+1}$ and condition
  on $(S_{t+1}, R_{t+1})$:
  \begin{equation*}
    q_\pi(s, a) = \sum_{s', r} p(s', r \mid s, a) \left( r + \gamma\,
      \E_\pi\left[G_{t+1} \given S_t = s, A_t = a, S_{t+1} = s',
        R_{t+1} = r\right] \right).
  \end{equation*}
  Given $S_{t+1} = s'$, the actions from $t+1$ on are drawn from $\pi$ and
  the transitions from $p$, independently of what happened before
  (Markov property \eqref{eq:markov} and the definition of a Markov
  policy). Hence the conditional expectation equals
  $\E_\pi[G_{t+1} \mid S_{t+1} = s'] = v_\pi(s')$.
\end{proof}

\begin{intuition}
  The value of a state is the average of the values of its actions,
  weighted by how often the policy chooses them. The value of an action is
  one step of reality (the reward and the next state, drawn by the
  environment) followed by the value of wherever the agent lands. $v$
  averages over the agent's choices, $q$ over the environment's.
\end{intuition}

\section{The Bellman expectation equations}

Substituting \eqref{eq:q-from-v} into \eqref{eq:v-from-q}, and conversely,
gives equations that involve one value function only.

\begin{theorem}{Bellman expectation equations}{bellman-exp}
  For every policy $\pi$,
  \begin{align}
    \label{eq:bellman-v}
    v_\pi(s) &= \sum_a \pi(a \mid s) \sum_{s', r} p(s', r \mid s, a)
      \left[ r + \gamma\, v_\pi(s') \right]
      \quad \text{for all } s, \\
    \label{eq:bellman-q}
    q_\pi(s, a) &= \sum_{s', r} p(s', r \mid s, a)
      \left[ r + \gamma \sum_{a'} \pi(a' \mid s')\, q_\pi(s', a') \right]
      \quad \text{for all } s, a.
  \end{align}
  Moreover $v_\pi$ is the \emph{unique} solution of \eqref{eq:bellman-v}
  and $q_\pi$ the unique solution of \eqref{eq:bellman-q}.
\end{theorem}

\begin{intuition}
  The Bellman equation is a consistency condition, not a recipe: the value
  of a state must equal what one expects to get in one step plus the
  discounted value of where one lands. It links the value of each state to
  the values of its neighbours, so that a reward at the end of a game is
  ``felt'' backwards, step by step, in all the states that lead to it.
  The equations form a linear system with one unknown per state, and
  \cref{prop:matrix} solves it.
\end{intuition}

The existence part is \cref{prop:v-q}. For uniqueness, and for everything
that follows, it pays to write the equations with operators.

\begin{definition}{Bellman operators}{operators}
  Let $\R^{\States}$ be the space of functions $v : \States \to \R$ with
  the sup norm $\norm{v}_\infty = \max_s \abs{v(s)}$. The \emph{Bellman
  expectation operator} of $\pi$ and the \emph{Bellman optimality operator}
  are the maps $\R^{\States} \to \R^{\States}$
  \begin{align}
    (\T^\pi v)(s) &= \sum_a \pi(a \mid s) \left[ r(s, a) + \gamma \sum_{s'}
      p(s' \mid s, a)\, v(s') \right], \\
    (\T^\star v)(s) &= \max_a \left[ r(s, a) + \gamma \sum_{s'}
      p(s' \mid s, a)\, v(s') \right].
  \end{align}
\end{definition}

With this notation, \eqref{eq:bellman-v} reads $v_\pi = \T^\pi v_\pi$:
$v_\pi$ is a fixed point of $\T^\pi$.

\begin{lemma}{Monotonicity and contraction}{contraction}
  For all $u, v \in \R^{\States}$:
  \begin{enumerate}[nosep, label=(\roman*)]
    \item if $u \leq v$ (pointwise) then $\T^\pi u \leq \T^\pi v$ and
      $\T^\star u \leq \T^\star v$;
    \item $\norm{\T^\pi u - \T^\pi v}_\infty \leq \gamma \norm{u -
      v}_\infty$ and $\norm{\T^\star u - \T^\star v}_\infty \leq \gamma
      \norm{u - v}_\infty$.
  \end{enumerate}
\end{lemma}

\begin{proof}
  (i) The coefficients $\pi(a \mid s)$ and $\gamma\, p(s' \mid s, a)$ are
  nonnegative, so each bracket is a nondecreasing function of $v$, and so
  are their weighted sum and their maximum.

  (ii) For $\T^\pi$: for every $s$,
  \begin{align*}
    \abs{(\T^\pi u)(s) - (\T^\pi v)(s)}
    &= \gamma \abs{\sum_a \pi(a \mid s) \sum_{s'} p(s' \mid s, a)
      \left(u(s') - v(s')\right)} \\
    &\leq \gamma \sum_a \pi(a \mid s) \sum_{s'} p(s' \mid s, a)
      \norm{u - v}_\infty,
  \end{align*}
  and the double sum of probabilities equals $1$. For $\T^\star$ we use the
  elementary inequality, valid for any functions $f, g$ on a finite set,
  \begin{equation}
    \label{eq:max-ineq}
    \abs{\max_a f(a) - \max_a g(a)} \leq \max_a \abs{f(a) - g(a)}.
  \end{equation}
  (Let $a_f$ maximise $f$; then $\max f - \max g \leq f(a_f) - g(a_f) \leq
  \max_a \abs{f(a) - g(a)}$, and symmetrically.) Applying it to
  $f(a) = r(s, a) + \gamma \sum_{s'} p(s' \mid s, a)\, u(s')$ and the
  corresponding $g$ with $v$, the rewards cancel and we conclude as for
  $\T^\pi$.
\end{proof}

\begin{intuition}
  Applying $\T^\pi$ to a guess $v$ means: look one step ahead with the true
  rewards and dynamics, and use the guess only for the states one step
  later. Whatever error the guess had, it now enters only through the next
  states, discounted by $\gamma$: each application pushes the error one
  step into the future, where it counts less. Monotonicity says that a
  better guess for the future gives a better estimate now. The inequality
  \eqref{eq:max-ineq} says that taking the best action cannot amplify
  errors: if every option is known to within $\delta$, the best option is
  also known to within $\delta$.
\end{intuition}

\begin{proof}[Proof of uniqueness in \cref{thm:bellman-exp}]
  If $v = \T^\pi v$ and $w = \T^\pi w$, then
  \begin{equation*}
    \norm{v - w}_\infty = \norm{\T^\pi v - \T^\pi w}_\infty
    \leq \gamma \norm{v - w}_\infty,
  \end{equation*}
  so $\norm{v - w}_\infty = 0$ because $\gamma < 1$. The equation for $q_\pi$
  is treated in the same way on $\R^{\States \times \Actions}$.
\end{proof}

\begin{proposition}{Matrix form}{matrix}
  Number the states, and let $r_\pi \in \R^{\States}$ and $P_\pi \in
  \R^{\States \times \States}$ be defined by $r_\pi(s) = \sum_a \pi(a \mid
  s)\, r(s, a)$ and $P_\pi(s, s') = \sum_a \pi(a \mid s)\, p(s' \mid s, a)$.
  Then $I - \gamma P_\pi$ is invertible and
  \begin{equation}
    v_\pi = (I - \gamma P_\pi)^{-1} r_\pi
    = \sum_{k=0}^{\infty} \gamma^k P_\pi^k\, r_\pi.
  \end{equation}
\end{proposition}

\begin{proof}
  The Bellman equation reads $v_\pi = r_\pi + \gamma P_\pi v_\pi$. The
  matrix $P_\pi$ is stochastic (nonnegative entries, rows summing to $1$),
  so its operator norm for $\norm{\cdot}_\infty$ is $1$ and $\norm{\gamma
  P_\pi}_\infty = \gamma < 1$. The Neumann series $\sum_k (\gamma P_\pi)^k$
  then converges and is the inverse of $I - \gamma P_\pi$: indeed
  $(I - \gamma P_\pi) \sum_{k=0}^{n} (\gamma P_\pi)^k = I - (\gamma
  P_\pi)^{n+1} \to I$. The term $\gamma^k P_\pi^k r_\pi$ is the expected
  discounted reward collected $k$ steps ahead.
\end{proof}

\begin{intuition}
  $P_\pi^k(s, s')$ is the probability of being in $s'$ after $k$ steps from
  $s$, so $P_\pi^k r_\pi$ is the reward expected at step $k$. The value is
  the sum of the rewards expected at each future step, discounted: the
  series $\sum_k \gamma^k P_\pi^k r_\pi$ is the definition of the value,
  and $(I - \gamma P_\pi)^{-1}$ is its closed form, the matrix analogue of
  $1 + x + x^2 + \cdots = 1/(1 - x)$.
\end{intuition}

\section{Optimal policies}

Value functions order policies: we say that $\pi \geq \pi'$ if $v_\pi(s)
\geq v_{\pi'}(s)$ for \emph{all} states $s$. It is not obvious that a policy
better than all the others in every state exists; the Bellman optimality
equation shows that it does.

\begin{definition}{Optimal value functions}{optimal}
  The \emph{optimal value functions} are
  \begin{equation}
    v_\star(s) = \max_\pi v_\pi(s), \qquad
    q_\star(s, a) = \max_\pi q_\pi(s, a),
  \end{equation}
  the maxima being over all stationary Markov policies. A policy $\pi_\star$
  is \emph{optimal} if $v_{\pi_\star} = v_\star$.
\end{definition}

\begin{theorem}{Bellman optimality equation}{bellman-opt}
  \begin{enumerate}[nosep, label=(\roman*)]
    \item The operator $\T^\star$ has a unique fixed point $v^\star$, and
      for every $v \in \R^{\States}$, $(\T^\star)^n v \to v^\star$.
    \item Let $\pi_\star$ be a \emph{greedy} policy with respect to
      $v^\star$: for every $s$, $\pi_\star(s)$ is an action maximising
      $r(s, a) + \gamma \sum_{s'} p(s' \mid s, a)\, v^\star(s')$. Then
      $v_{\pi_\star} = v^\star$ and $v_\pi \leq v^\star$ for every policy
      $\pi$. Hence $v_\star = v^\star$, and a deterministic optimal policy
      exists.
    \item The optimal action values are $q_\star(s, a) = r(s, a) + \gamma
      \sum_{s'} p(s' \mid s, a)\, v_\star(s')$, and they satisfy
      \begin{equation}
        \label{eq:bellman-q-opt}
        v_\star(s) = \max_a q_\star(s, a), \qquad
        q_\star(s, a) = \sum_{s', r} p(s', r \mid s, a)
          \left[ r + \gamma \max_{a'} q_\star(s', a') \right].
      \end{equation}
  \end{enumerate}
\end{theorem}

\begin{proof}
  (i) $\R^{\States}$ with $\norm{\cdot}_\infty$ is complete and $\T^\star$
  is a $\gamma$-contraction (\cref{lem:contraction}); the Banach fixed point
  theorem gives the unique fixed point and the convergence of the
  iterates. (We recall its proof: $\norm{(\T^\star)^{n+1} v - (\T^\star)^n
  v}_\infty \leq \gamma^n \norm{\T^\star v - v}_\infty$, so the iterates
  form a Cauchy sequence; its limit is a fixed point by continuity of
  $\T^\star$, unique by the argument used for \cref{thm:bellman-exp}.)

  (ii) By definition of the greedy policy, $\T^{\pi_\star} v^\star =
  \T^\star v^\star = v^\star$: $v^\star$ is a fixed point of
  $\T^{\pi_\star}$, whose unique fixed point is $v_{\pi_\star}$
  (\cref{thm:bellman-exp}); so $v_{\pi_\star} = v^\star$. Now let $\pi$ be
  any policy. For every $v$, $\T^\pi v \leq \T^\star v$, because a weighted
  average of the brackets is at most their maximum. Hence $\T^\pi v^\star
  \leq \T^\star v^\star = v^\star$, and by monotonicity
  (\cref{lem:contraction}), by induction, $(\T^\pi)^n v^\star \leq
  v^\star$ for all $n$: if $(\T^\pi)^n v^\star \leq v^\star$ then
  $(\T^\pi)^{n+1} v^\star \leq \T^\pi v^\star \leq v^\star$. Since
  $(\T^\pi)^n v^\star \to v_\pi$ (Banach again, for $\T^\pi$), we get
  $v_\pi \leq v^\star$.

  (iii) By \eqref{eq:q-from-v} applied to $\pi_\star$, the right-hand side
  is $q_{\pi_\star}(s, a)$. For any policy $\pi$, \eqref{eq:q-from-v} and
  $v_\pi \leq v_\star$ give $q_\pi(s, a) \leq q_{\pi_\star}(s, a)$; so
  $q_\star = q_{\pi_\star}$. The fixed point equation $v_\star = \T^\star
  v_\star$ is the first identity of \eqref{eq:bellman-q-opt}; substituting
  it into the definition of $q_\star$ gives the second.
\end{proof}

\begin{intuition}
  This is Bellman's \emph{principle of optimality}: whatever the first
  move, an optimal plan must continue optimally from wherever it lands, so
  the value of playing well is the best, over the first move, of the reward
  plus the value of playing well afterwards. $q_\star(s, a)$ answers the
  question ``what if I play $a$ now and play perfectly afterwards?'', and
  the best action is the one with the best answer. Each application of
  $\T^\star$ makes the agent look one move further ahead, and the
  contraction says that, far enough ahead, the starting guess no longer
  matters.
\end{intuition}

\begin{remark}{What the theorem says}{opt-meaning}
  An optimal policy needs no randomness and no memory: in each state it is
  enough to play an action maximising $q_\star(s, \cdot)$. Knowing
  $q_\star$ is therefore enough to act optimally, \emph{without} knowing
  the dynamics $p$; this is the starting point of Q-learning and DQN
  (\cref{ch:dqn}). Knowing $v_\star$ only is not enough to act: choosing
  the greedy action requires $p$, i.e.\ a model of the environment. (One
  can also show that policies depending on the whole history cannot do
  better than $v_\star$; see Puterman, \emph{Markov Decision Processes},
  1994, Theorem 6.2.10.)
\end{remark}

\begin{pacman}{To the power cookie, or away?}{power-or-flee}
  Pacman is at a junction $s$, with a ghost two cells away. Consider two
  options, with simplified consequences.
  \begin{itemize}[nosep]
    \item $a_1$, \emph{run to the power cookie}: with probability $1 - \rho$
      Pacman gets there first, receives about $5$ and reaches a state worth
      $25$ (the ghost will be eaten); with probability $\rho$ the ghost
      catches him ($-35$, game over).
    \item $a_2$, \emph{flee down the corridor}: Pacman pays the living cost
      $-0.15$; with probability $0.6$ the ghost follows and Pacman is back
      in a situation like $s$; with probability $0.4$ he escapes to a state
      worth $10$.
  \end{itemize}
  With $\gamma = 0.99$, \eqref{eq:bellman-q-opt} gives
  \begin{equation*}
    q_\star(s, a_1) = (1 - \rho)(5 + 0.99 \times 25) - 35\,\rho, \qquad
    q_\star(s, a_2) = -0.15 + 0.99 \left( 0.6\, v_\star(s) + 0.4 \times 10
    \right),
  \end{equation*}
  with $v_\star(s) = \max(q_\star(s, a_1), q_\star(s, a_2))$: the value of
  fleeing depends on the value of the state itself.
  \begin{itemize}[nosep]
    \item For $\rho = 0.3$: $q_\star(s, a_1) = 10.325$. Assuming $v_\star(s)
      = 10.325$ gives $q_\star(s, a_2) = 9.943 < 10.325$, which is
      consistent: run to the power cookie, $v_\star(s) = 10.325$.
    \item For $\rho = 0.4$: $q_\star(s, a_1) = 3.85$. If fleeing is optimal,
      $v_\star(s) = q_\star(s, a_2)$ solves $q = -0.15 + 0.99 (0.6\, q + 4)$,
      i.e.\ $q = 3.81 / 0.406 = 9.384 > 3.85$, which is consistent: flee.
  \end{itemize}
  A $10\%$ change in the risk changes the optimal decision. Learning such
  trade-offs from experience is what the agents of the repository do.
\end{pacman}

\section{Other criteria}

The discounted return is not the only possible objective.
\begin{itemize}
  \item \textbf{Finite horizon.} With a fixed number $H$ of steps and no
    discount, the optimal value depends on the number of steps left: $v_h$
    for $h$ steps to go satisfies $v_0 = 0$ and $v_{h} = \T^\star v_{h-1}$
    (with $\gamma = 1$), and the optimal policy is non-stationary. This is
    backward induction, the original form of dynamic programming
    (Bellman, 1957).
  \item \textbf{Undiscounted episodic tasks.} With $\gamma = 1$, the
    returns are finite if every policy reaches a terminal state with
    probability $1$ (\emph{proper} policies); the results of this chapter
    extend to this \emph{stochastic shortest path} setting (Bertsekas and
    Tsitsiklis, 1996).
  \item \textbf{Average reward.} For continuing tasks one can maximise
    $\lim_{T \to \infty} \frac{1}{T} \E[\sum_{t=1}^T R_t]$. The values are
    then defined relative to the average reward $\bar r_\pi$, $v_\pi(s) =
    \E_\pi[\sum_t (R_{t+1} - \bar r_\pi) \mid S_0 = s]$, and satisfy
    $v_\pi(s) = \sum_a \pi(a \mid s) \sum_{s', r} p(s', r \mid s, a) [r -
    \bar r_\pi + v_\pi(s')]$. Discounting with $\gamma$ close to $1$ is a
    common approximation of this criterion.
\end{itemize}

\section{Summary}

\begin{itemize}[nosep]
  \item An MDP is described by states, actions, a transition-reward kernel
    $p$ satisfying the Markov property, and a discount factor $\gamma$.
  \item The value functions $v_\pi$ and $q_\pi$ are expected returns; they
    are the unique fixed points of the Bellman expectation operators.
  \item The optimal value $v_\star$ is the unique fixed point of the
    Bellman optimality operator, a $\gamma$-contraction; any greedy policy
    with respect to $v_\star$ (or $q_\star$) is optimal.
  \item For Pacman, the state is the board configuration; the feature
    planes of game \code{v1} contain it entirely, while game \code{v2} is
    partially observable.
\end{itemize}
The next chapters turn these equations into algorithms: first when $p$ is
known (dynamic programming), then when it must be replaced by experience.
